% TOP_PROBLEM: {"problemId": "problem.top500-omission-w045049-f-conjecture-mbar-zero-n", "canonicalTitle": "F-Conjecture for the Nef Cone of Mbar_(0,n)", "releaseRank": 499, "primaryDomain": "geometry_topology", "primaryDomainLabel": "Geometry & topology", "displayStatus": "open_with_solved_subcases", "releaseStatus": "open_with_solved_subcases"}
% !TeX program = lualatex
% ATTEMPT_STATUS: unresolved
% Model/process: GPT_6_Astra_Ultra; autonomous mathematical attempt,
% primary-source audit, exact finite diagnostics, and adversarial review.
% Prepared 27 September 2026; no general solution or novelty claimed.
\documentclass[11pt]{article}
\usepackage[margin=25mm]{geometry}
\usepackage{fontspec}
\setmainfont{Latin Modern Roman}
\usepackage{amsmath,amssymb,mathtools}
\usepackage{unicode-math}
\setmathfont{Latin Modern Math}
\usepackage{xurl,hyperref}
\hypersetup{colorlinks=true,urlcolor=blue}
\setcounter{secnumdepth}{0}
\setlength{\parindent}{0pt}
\setlength{\parskip}{5pt}
\setlength{\emergencystretch}{3em}
\begin{document}
\section*{\texorpdfstring{F-Conjecture for the Nef Cone of Mbar\_(0,n)}{F-Conjecture for the Nef Cone of Mbar\_(0,n)}}\label{499}
\subsection{Definitions and mathematical statement}
The projective moduli space $\overline M_{0,n}$ parametrizes stable genus-zero nodal curves with $n$ distinct labeled smooth marked points. Stability requires at least three marked points or nodes on each irreducible component. An F-curve is the one-dimensional family obtained by partitioning the labels into four nonempty groups, attaching the corresponding fixed tails, and varying the cross-ratio on the central four-pointed component, including its stable limits.

For a real numerical divisor class $D$, nefness means $D\cdot C\ge0$ for every complete curve $C$. The conjecture says
\[
 D\text{ nef}\quad\Longleftrightarrow\quad
 D\cdot F\ge0\text{ for every F-curve }F,
\]
or dually that $\overline{\operatorname{NE}}_1(\overline M_{0,n})$ is generated by their classes. This is for every $n\ge3$. For $n=3$ the space is a point; for $n=4$ the relevant F-curve is $\overline M_{0,4}$ itself. Thus the phrase boundary stratum uses the conventional extension to this case, not a requirement that the curve lie in the proper boundary.
\par\smallskip\textit{Formulation and concept references:} \href{https://arxiv.org/html/2507.12434v1}{[S1]}.

\subsection{Short English statement}
Can nef divisors on every genus-zero pointed-curve moduli space be recognized by testing only the standard one-dimensional F-curves?
\subsection{Research attempt}

\paragraph{Conventions and source audit (27 September 2026).}
The calculation is over the complex numbers, with real numerical divisor
and curve classes. It addresses the usual characteristic-zero formulation;
no argument about reduction in positive characteristic is asserted. Write
$X_n=\overline M_{0,n}$ and call a divisor \emph{F-nef} if it has nonnegative
intersection with every F-curve. There are finitely many partitions of
$[n]$ into four nonempty blocks, hence their classes generate a closed
polyhedral cone. The easy implication is nef $\Rightarrow$ F-nef. The task
is the reverse implication for every $n$, including divisors on faces of
the cone; strict inequalities cannot be substituted for weak ones.

The supplied source [S1], version 1 of 16 July 2025, states equivalence
between the symmetric and unrestricted conjectures when each is quantified
over all numbers of markings. Its Theorem 4.2 realizes an F-nef divisor as
a pullback of a symmetric F-nef divisor after attaching tails and increasing
the number of markings. Its Theorem 1.3 proves the strong F-conjecture for
$X_8$. These are the manuscript's stated results; the present note does not
reprove or independently certify them. In particular, ordinary averaging
at a fixed $n$ is not its reduction. The source was directly inspected,
rather than inferred from its title.

The attack below first solves the elementary $n=5$ cone problem by plane
curve intersection and an explicit positive decomposition. It then attempts
an induction through the boundary, retaining the exact missing assertion
for curves which meet the interior. The known low-dimensional result is
used as a diagnostic for the method, not represented as a new advance.

\paragraph{Intersection bookkeeping before attempting an induction.}
For a stable unordered partition $I\mid I^c$ let $\Delta_I=\Delta_{I^c}$
be the corresponding boundary divisor, and denote the cotangent classes by
$\psi_i$. If $A,B,C,D$ are the four blocks defining $F$, then
\begin{equation}\label{ra499:intersections}
 \psi_i\cdot F=\begin{cases}1,&\{i\}\text{ is a block},\\0,&\text{otherwise},\end{cases}
 \qquad
 \Delta_I\cdot F=\begin{cases}
 1,&I\mid I^c\text{ is a union of two blocks versus the other two},\\
 -1,&I\text{ or }I^c\text{ is a nonsingleton block},\\
 0,&\text{otherwise}.
 \end{cases}
\end{equation}
Each unordered partition is counted once. In particular there are three,
not six, positive terms in the boundary formula.

Here is a geometric explanation of the signs. The varying central
four-pointed component meets three boundary points of $X_4\cong\mathbb P^1$;
they correspond to its three two-versus-two splittings and give transverse
intersection $+1$. A fixed tail contributes a boundary divisor containing
the family. The normal line to this gluing divisor is the tensor product
of the tangent lines at the two branches of its node. Its degree is the
negative cotangent degree on the moving central component, namely $-1$;
the tail cotangent degree is zero because that tail is fixed. At a marked
point on a fixed tail the cotangent line is constant, whereas a singleton
mark on the central component has cotangent degree one on $X_4$.
Choosing smooth fixed tails makes these the only relevant boundaries.
The numerical class is unchanged on specializing the tails, which recovers
the standard maximally degenerate representative. These are the usual
F-curve intersection conventions underlying [S1, Section 2].

Consequently, for any representation
\[
 L=\sum_i a_i\psi_i-\sum_{I\mid I^c}b_I\Delta_I,
\]
one obtains the concrete test
\[
 L\cdot F_{A,B,C,D}=
 \sum_{\{i\}\text{ a block}}a_i+
 \sum_{U\in\{A,B,C,D\},\ |U|\ge2}b_U
 -b_{A\cup B}-b_{A\cup C}-b_{A\cup D}.
\]
The coefficient representation need not be unique. This formula computes
intersection with the numerical class and does not require uniqueness.
An incorrectly reversed normal-bundle sign would alter the cone and
invalidate every later linear-programming test.

\paragraph{A complete small-dimensional calculation.}
For $n=3$ there are no curve classes. For $n=4$, $X_4\cong\mathbb P^1$,
and its single curve class is the F-curve itself. For $n=5$ use the
Kapranov realization [E1, Section 4.3]
\[
 X_5\cong\operatorname{Bl}_{p_1,p_2,p_3,p_4}\mathbb P^2,
\]
where no three of the four points are collinear. Let $H$ be the pullback
of a line and $E_i$ the exceptional curves. The ten boundary curves, which
are exactly the ten F-curves in this dimension, are
\[
 E_i\quad(1\le i\le4),\qquad
 L_{ij}=H-E_i-E_j\quad(1\le i<j\le4).
\]
The surface intersection pairing is
$H^2=1$, $H\cdot E_i=0$, and $E_i\cdot E_j=-\delta_{ij}$.
The blow-up description and the identification of its ten boundary curves
are the geometric input; the cone argument that follows is supplied in
full.

\textbf{Lemma 499.1.}
A real divisor $L=aH-\sum_{i=1}^4b_iE_i$ is nef if and only if
\begin{equation}\label{ra499:five}
 b_i\ge0\quad\text{for all }i,\qquad
 a-b_i-b_j\ge0\quad\text{for all }i<j.
\end{equation}

\emph{Proof.}
Necessity follows by intersecting $E_i$ and $L_{ij}$. For sufficiency let
$C$ be an irreducible complete curve. If $C$ is one of those ten curves,
its intersection is already nonnegative. Otherwise its plane image is an
irreducible curve of degree $d\ge1$, with multiplicities $m_i\ge0$ at the
four points, and
$[C]=dH-\sum m_iE_i$. This image is not any line $p_ip_j$.
Bezout's theorem with that line therefore gives
\begin{equation}\label{ra499:bezout}
 m_i+m_j\le d\quad(i<j).
\end{equation}
In particular $m_i\le d$ and
$\sum m_i\le2d$, the latter by adding the inequalities for any complementary
pair of pairs.

Relabel the coefficients so $b_1\ge b_2\ge b_3\ge b_4\ge0$, relabeling
the corresponding multiplicities at the same time. Then
\[
 \sum_i b_im_i
 \le b_1m_1+b_2\sum_{i=2}^4m_i
 \le b_1m_1+b_2(2d-m_1)
 \le d(b_1+b_2)\le ad.
\]
The penultimate inequality uses both $b_1-b_2\ge0$ and $m_1\le d$.
Thus $L\cdot C=ad-\sum b_im_i\ge0$ for every irreducible curve.
Linearity handles effective cycles, and continuity handles their numerical
closure. This proves nefness under precisely the weak inequalities in
\eqref{ra499:five}.\hfill$\square$

This proves the $n=5$ instance of the target using no assertion about
effective cone generators beyond the ten explicitly tested curves. It
also supplies the necessary closed-cone boundary cases such as $b_i=0$.

\paragraph{A constructive dual certificate for the same surface.}
The preceding proof has a useful positive-decomposition counterpart.
Suppose $d>0$ and numbers $m_i\ge0$ satisfy
\[
 m_i\le d,\qquad \sum_{i=1}^4m_i\le2d.
\]
Choose numbers $r_i$ with $m_i\le r_i\le d$ and $\sum r_i=2d$.
They exist by distributing the deficit $2d-\sum m_i$ among the four
capacities $d-m_i$: their sum is $4d-\sum m_i$, which is at least the
deficit. This is an explicit finite greedy filling procedure, with no
integrality requirement.

Partition a circle of circumference $2d$ into four consecutive half-open
arcs $I_i$ of lengths $r_i$. Pair each point $x$ with its antipode $x+d$.
No pair has both endpoints in the same arc, except possibly at endpoints
of measure zero, because $r_i\le d$. For $i<j$ let $x_{ij}$ be the
Lebesgue measure of the set of antipodal pairs joining $I_i$ to $I_j$,
where each pair is counted once using its representative $x\in[0,d)$.
Then
\[
 x_{ij}\ge0,\qquad \sum_{i<j}x_{ij}=d,
 \qquad \sum_{j\ne i}x_{ij}=r_i.
\]
The last equality counts the total arc length incident to pairs at vertex
$i$; every point of $I_i$ supplies exactly one incident endpoint.
It follows that
\begin{equation}\label{ra499:decomposition}
 dH-\sum_i m_iE_i
 =\sum_{i<j}x_{ij}L_{ij}+\sum_i(r_i-m_i)E_i.
\end{equation}
Every coefficient is nonnegative. Applying this to the multiplicities
obtained in \eqref{ra499:bezout}, and treating the ten excluded curves
separately, proves directly that the entire curve cone of $X_5$ is generated
by its F-curves. The decomposition may be real; that is exactly sufficient
for the real numerical cone in the question.

For example a conic through all four points has class
$2H-E_1-E_2-E_3-E_4=L_{12}+L_{34}$.
A general line has class $H=L_{12}+E_1+E_2$.
These are numerical decompositions, not assertions that every actual
smooth curve is set-theoretically a union of boundary components.

\paragraph{The induction that almost closes, and its precise extra input.}
Consider a boundary divisor
\[
 j:\Delta_I\cong X_{|I|+1}\times X_{n-|I|+1}\longrightarrow X_n.
\]
Both factors have fewer than $n$ markings. A divisor restriction splits
numerically as $j^*L=p_1^*L_1+p_2^*L_2$. One way to check this specific
splitting is to use the boundary-divisor generators of the Picard group:
a different boundary divisor restricts to boundary divisors on the relevant
factor, or is disjoint; the self-restriction is the negative sum of the
two cotangent classes at the gluing marks. Each summand is a pullback from
a factor. The generation and gluing conventions are those of [S1,
Section 2]; this does not invoke a general claim about Picard groups of
arbitrary products.

If $L$ is F-nef, each $L_i$ is F-nef. Indeed an F-curve in one factor
with the other factor fixed maps to an F-curve class in $X_n$: the gluing
mark on its variable component is replaced by the fixed tail containing
the other labels. The block containing that gluing mark is thereby
enlarged, and the other three nonempty blocks are unchanged. Its degree
under the pullback equals the degree under $L$. This also explains why
the statement concerns numerical classes, regardless of the chosen fixed
tail.

Assume inductively that the F-conjecture holds for every smaller number
of markings. Then $L_1,L_2$ are nef. Their pullbacks and their sum are
nef, so $L$ has nonnegative intersection with every curve contained in
every boundary divisor. This is a valid consequence of the induction
hypothesis. It leaves curves whose generic point is in $M_{0,n}$.
An irreducible curve meeting the interior need not be an F-curve merely
because its stable limits meet some boundary divisors.

\textbf{Lemma 499.2 (a sufficient extension step).}
Under the same lower-$n$ induction hypothesis, an F-nef divisor $L$ which
is numerically equivalent to an effective real combination of boundary
divisors is nef.

\emph{Proof.}
Write $L\equiv\sum c_I\Delta_I$ with $c_I\ge0$. A curve contained in
the boundary was handled by the restriction argument. For an irreducible
curve $C$ not contained in any boundary divisor, pull back the defining
section of $\Delta_I$ to the normalization of $C$. It is nonzero and has
an effective zero divisor, so $\Delta_I\cdot C\ge0$. Therefore
$L\cdot C=\sum c_I(\Delta_I\cdot C)\ge0$. These two cases cover every
irreducible complete curve.\hfill$\square$

The missing step is consequently very concrete: find such a nonnegative
boundary representation for every F-nef $L$, or obtain a different bound
on its degree on every curve meeting the interior. F-inequalities are
necessary linear inequalities for that representation but the argument
has not shown them sufficient. Substituting the desired representation
without proof would be the first circular step. The strong statement in
[S1, Theorem 1.3] gives it in a specified small case, not for all $n$.

\paragraph{Two attempted shortcuts and their stress tests.}
First, a finite inequality system computes the dual of the cone generated
by the listed F-curves. This equals the desired nef cone only after one
proves the geometric assertion about all curves. Linear feasibility alone
cannot establish that assertion. The surface calculation succeeds because
Bezout gives extra inequalities for an arbitrary irreducible curve; a
direct replacement with such force in all dimensions is not supplied.

Second, averaging $L$ over permutations proves only a statement about
$\overline L=(n!)^{-1}\sum_\sigma\sigma^*L$. Averaging preserves F-nefness
and nefness, but nefness of the average need not imply nefness of its
summands. A complete elementary model of this logical failure is
\[
 K=\{(x,y):x\ge0,\ y\ge0\},\qquad
 F=\{(x,y):x+y\ge0\},
\]
with the group exchanging coordinates. The invariant slices of $K$ and
$F$ agree, but $(-1,2)\in F\setminus K$, while its average
$(1/2,1/2)$ is in $K$. This is a counterexample to the abstract averaging
inference, not a divisor counterexample to the F-conjecture. The attaching
construction in [S1] is substantially stronger than this averaging step.

The finite verification record enumerates all 13,446 integer multiplicity tuples
$1\le d\le12$, $0\le m_i\le d$ obeying the six Bezout inequalities.
For every tuple it constructs the antipodal certificate by exact rational
interval arithmetic, verifies all coefficients and all five coordinates
in \eqref{ra499:decomposition}, and records the total tested. It does not
assert that each numerical tuple is realized by an irreducible plane
curve; the certificate applies to the larger numerical test set. No
floating-point feasibility tolerance or inferred all-$n$ conclusion is used.

\paragraph{Outcome.}
\textbf{Outcome: UNRESOLVED, with a complete known small-case reconstruction.}
The work supplies a full $n\le5$ proof, an explicit positive decomposition
algorithm for the surface case, and a proved conditional induction step.
The first unproved assertion is control of F-nef divisors on curves which
meet the interior, for example through an effective boundary representation
in arbitrary $n$. No such representation is constructed in general and no
F-nef non-nef divisor is found. The general catalogue statement is neither
proved nor disproved.

\subsection{Sources}
\begin{itemize}
\item[S1]M. Fedorchuk and A. Mellit, Symmetric and non-symmetric F-conjectures are equivalent, arXiv 2507.12434.
\url{https://arxiv.org/html/2507.12434v1}.
\textit{Formulation source.}
\item[E1]M. M. Kapranov, \emph{Chow quotients of Grassmannians I},
arXiv:alg-geom/9210002, Section 4.3, especially Theorem 4.3.3;
the case $n=5$ gives the blow-up of four points in the plane.
\url{https://arxiv.org/pdf/alg-geom/9210002}.
\textit{Primary manuscript inspected 27 September 2026.}
\end{itemize}
\end{document}
